\documentclass[11pt]{article}
\usepackage[a4paper,margin=1in]{geometry}
\usepackage{amsmath,amssymb,amsthm}
\usepackage{booktabs}
\usepackage{xcolor}
\usepackage{hyperref}
\hypersetup{colorlinks=true,linkcolor=blue,urlcolor=blue,citecolor=blue}
\usepackage{enumitem}

\newcommand{\fn}[1]{\texttt{#1}}
\newcommand{\Tr}{\operatorname{Tr}}
\newcommand{\Phid}{\Phi^{\dagger}}
\newcommand{\B}{\mathcal{B}}
\newcommand{\A}{\mathcal{A}}
\newcommand{\C}{\mathcal{C}}
\newcommand{\F}{\mathcal{F}}
\newcommand{\D}{\mathcal{D}}
\newcommand{\HH}{\mathcal{H}}
\newcommand{\M}{\mathcal{M}}
\newcommand{\term}[1]{\medskip\noindent\textbf{#1}\quad}

\theoremstyle{plain}
\newtheorem{theorem}{Theorem}
\newtheorem{proposition}{Proposition}
\theoremstyle{remark}
\newtheorem{remark}{Remark}

\title{\textbf{Report R25 --- the algebraic treatment}\\
\large A reading path and glossary for a physicist who learned quantum
mechanics the Hilbert-space way}
\author{Tier 1 analysis}
\date{2026-08-22}

\begin{document}
\sloppy
\emergencystretch=3em
\maketitle

\begin{abstract}
Reports R18--R24 of this project are written in a language --- commutants, von
Neumann algebras, multiplicative domains, noiseless subsystems --- that a
standard physics curriculum does not teach. This document is the missing
prerequisite, assembled for thesis use. \S\ref{sec:pivot} states the five
conceptual moves that separate the algebraic formulation from the one taught in
a physics degree, since knowing \emph{which} habits to unlearn is most of the
difficulty. \S\ref{sec:path} gives a reading path with an explicit minimum:
three items suffice to work in this area, and the rest is reference material.
\S\ref{sec:theorems} states the three structure theorems everything else rests
on, in a form a physicist can use. \S\ref{sec:glossary} is a glossary of some
forty terms, each with a definition, an intuition, and the reference to cite.
\S\ref{sec:collisions} lists the notational collisions in this literature that
are known to cause errors --- including three distinct uses of ``strong versus
weak'', two of which this project uses. \S\ref{sec:project} maps the vocabulary
onto where each piece is used in R18--R24, so a thesis chapter can cite in the
right direction.

Every journal reference marked $\checkmark$ in \S\ref{sec:biblio} was verified
against the publisher or arXiv while this document was written. Textbooks are
marked $\dagger$: their bibliographic data is standard and was not
re-verified, so check the edition before citing a page number.
\end{abstract}

\tableofcontents

\section{Five moves a physics degree does not teach}
\label{sec:pivot}

The algebraic formulation is not a different theory. It is the same quantum
mechanics with the logical order reversed, and the reversal buys structure
theorems that are invisible in the usual presentation. Five specific habits
have to be given up.

\paragraph{1. Observables are primary; states are functionals on them.}
The usual order is: fix a Hilbert space $\HH$; states are unit vectors or
density matrices; observables are self-adjoint operators. The algebraic order
is: fix a $C^*$-algebra $\A$ of observables, and a \emph{state} is a positive
linear functional $\omega : \A \to \mathbb{C}$ with $\omega(\mathbb{1}) = 1$.
The Hilbert space is then \emph{derived} from the algebra and the state, by the
GNS construction. For finite-dimensional problems this reversal changes nothing
computationally --- but it is what makes ``the algebra of observables that
survive'' a well-posed object, and that object is the subject of R22 and R24.

\paragraph{2. A subsystem need not be a tensor factor of the whole.}
This is the most important move, and the one with no counterpart in an
undergraduate course. In a physics degree, ``subsystem'' means
$\HH = \HH_A \otimes \HH_B$ globally. The structure theorem
(\S\ref{sec:wedderburn}) says the general situation is
\[
  \HH \;\cong\; \bigoplus_k \; \HH_{n_k} \otimes \HH_{m_k},
\]
a tensor product \emph{inside each block} of a direct sum. A subsystem in the
algebraic sense is one of the $\HH_{n_k}$ factors: a genuine quantum degree of
freedom, living inside a superselection block, which does not factor out of
$\HH$ globally. Every decoherence-free subsystem, every noiseless subsystem and
every subsystem code \cite{KLP2005} is of this form. If only one of the two
readings of ``subsystem'' is in your head, this literature is confusing at
every step.

\paragraph{3. A symmetry is an algebra, not a group.}
The usual object is a group $G$ acting by unitaries $U_g$. The algebraic object
is the \emph{commutant} $\{U_g\}' = \{A : [A, U_g] = 0\ \forall g\}$, and it is
the commutant that the structure theorems see. Two consequences. A group is
only one way of generating a commutant, and the interesting commutants in
fragmentation are generated by no group at all --- that is precisely Moudgalya
and Motrunich's point \cite{MM2022}. And degeneracy is a statement about the
commutant being non-Abelian, not about the group being non-Abelian.

\paragraph{4. Conserved quantities need not close under multiplication.}
In closed dynamics, $A$ conserved means $[A,H]=0$, and the conserved quantities
form an algebra. That is an accident of reversibility. For an irreversible
channel $\Phi$ the conserved quantities are $\F = \{A : \Phid(A) = A\}$, a
self-adjoint unital \emph{operator system} but in general not an algebra:
$A, B \in \F$ does not give $AB \in \F$. It becomes an algebra exactly when
$\Phi$ has a faithful invariant state (Frigerio \cite{Frigerio1978}). The
failure is physical, not a technicality --- R22 and R24 turn on it.

\paragraph{5. Six fields have independently found one theorem.}
Superselection sectors, decoherence-free subspaces, noiseless subsystems,
operator-algebra error correction, information-preserving structures, and
Krylov/commutant fragmentation are, mathematically, one structure theorem seen
from six directions. This is good news --- results transfer --- and bad news,
because each field kept its own notation. \S\ref{sec:collisions} lists the
traps.

\section{A reading path}
\label{sec:path}

\subsection{If you read only three things}

For working in the area of this project, these three suffice; the rest of
\S\ref{sec:path} is reference material.

\begin{enumerate}[label=\textbf{\arabic*.}]
\item \textbf{Wolf, \emph{Quantum Channels \& Operations: Guided Tour}}
\cite{Wolf2012}. Free lecture notes, finite-dimensional throughout, and the
best single source for the structure theory of channels: Kraus, Stinespring,
Choi, the multiplicative domain, fixed points, peripheral spectrum. Written for
people who want to \emph{use} it. Start here.

\item \textbf{Albert, \emph{Lindbladians with multiple steady states}}
\cite{Albert2018}. A 134-page thesis that reads as a review, and the most
physicist-friendly account of the asymptotic decomposition
$\rho(\infty) = \bigoplus_k p_k\, \rho_k \otimes \sigma_k$ on which R22 and R24
are built. Chapters 1--3 are the relevant ones.

\item \textbf{Moudgalya \& Motrunich, \emph{Hilbert Space Fragmentation and
Commutant Algebras}} \cite{MM2022}. The paper that made ``commutant algebra''
the standard language for fragmentation, and whose criterion this project's
$n_{\mathrm{wcc}}$ literally computes.
\end{enumerate}

\noindent Add \textbf{Moudgalya, Bernevig \& Regnault} \cite{MBR2022} if a
thesis introduction needs a broad review of fragmentation and scars to cite; it
is the standard one and is pedagogical.

\subsection{Foundations: what an algebraic formulation is}

\begin{itemize}
\item \textbf{Strocchi} \cite{Strocchi2008} --- short, physicist-facing, the
gentlest account of \emph{why} one would formulate quantum mechanics
algebraically. Read before any mathematics book.
\item \textbf{Landsman} \cite{Landsman2017} --- open access, comprehensive,
with self-contained appendices on functional analysis and $C^*$-algebras. The
best single reference for the foundations proper. Use as a reference, not a
cover-to-cover read.
\item \textbf{Bratteli \& Robinson}, vol.~1 \cite{BR1} --- the standard
mathematical-physics source. Volume~1 covers $C^*$- and $W^*$-algebras,
symmetry groups and decomposition of states, which is the relevant part;
volume~2 is equilibrium statistical mechanics and is not needed here.
\item \textbf{Haag} \cite{Haag1996} --- where the superselection language comes
from. Cite for provenance; do not read for technique.
\end{itemize}

\noindent For the pure mathematics, if a proof is ever needed:
\textbf{Murphy} \cite{Murphy1990} is the readable introduction,
\textbf{Blackadar} \cite{Blackadar2006} the encyclopaedia, \textbf{Takesaki}
\cite{Takesaki2002} the definitive treatment. In finite dimensions --- where
this entire project lives --- none is strictly necessary: the structure theorem
is Wedderburn--Artin, and it is linear algebra.

\subsection{Quantum channels and open-system structure}

\begin{itemize}
\item \textbf{Wolf} \cite{Wolf2012} --- the primary source, as above.
\item \textbf{Watrous} \cite{Watrous2018} --- rigorous, finite-dimensional;
best for precise statements of Choi and Stinespring, and for anything
information-theoretic.
\item \textbf{Heinosaari \& Ziman} \cite{HZ2011} --- states, effects and
channels done carefully; useful if measurement theory is also needed.
\item \textbf{Breuer \& Petruccione} \cite{BP2002}, \textbf{Rivas \& Huelga}
\cite{RH2012} --- the standard open-systems texts. Use for the physics of
Lindblad dynamics; neither emphasises algebraic structure.
\item \textbf{Baumgartner \& Narnhofer} \cite{BN2008} --- the structure of the
set of stationary states, and the source of the term \emph{enclosure}.
\item \textbf{Carbone \& Pautrat} \cite{CP2016} --- the irreducible
decomposition done carefully, with the transient/recurrent split explicit. The
most mathematically precise of the group.
\item \textbf{Ticozzi \& Viola} \cite{TV2008} --- invariance versus
\emph{attractivity}. Worth knowing: the gap between an invariant subspace and
one actually reached is exactly what R24 flags between its Prop.~3 and its
absorption measurements.
\end{itemize}

\subsection{Symmetry in open systems}

\begin{itemize}
\item \textbf{Bu\v{c}a \& Prosen} \cite{Buca2012} --- the strong/weak symmetry
distinction. Short, and the definition every later paper uses.
\item \textbf{Albert \& Jiang} \cite{AJ2014} --- conserved quantities and
asymptotic structure, stated for physicists.
\item \textbf{Lieu \emph{et al.}} \cite{Lieu2020} --- what symmetry breaking
means for a Lindbladian, and the observation that a strong-symmetry-broken
phase can protect a qubit where a weak one protects only a classical bit.
Directly relevant to R22 \S8.
\item \textbf{Lessa \emph{et al.}} \cite{Lessa2024} --- strong-to-weak
spontaneous symmetry breaking in mixed states; a third use of ``strong/weak''
(\S\ref{sec:collisions}).
\end{itemize}

\subsection{Fragmentation and commutant algebras}

\begin{itemize}
\item \textbf{Sala \emph{et al.}} \cite{Sala2020} and \textbf{Khemani, Hermele
\& Nandkishore} \cite{Khemani2020} --- the two papers that introduced
fragmentation, independently and simultaneously. Cite both.
\item \textbf{Moudgalya \& Motrunich} \cite{MM2022} --- the commutant-algebra
formulation. Key definitions: the \emph{bond algebra} generated by the local
terms, its \emph{commutant}, and the criterion that fragmentation is
exponential growth of the commutant's dimension.
\item \textbf{Moudgalya \& Motrunich} \cite{MM2024} --- symmetries as
frustration-free ground states of a local superoperator; a computational route
to commutants.
\item \textbf{Moudgalya, Bernevig \& Regnault} \cite{MBR2022} --- the review.
\item \textbf{Li, Sala \& Pollmann} \cite{Li2023} --- fragmentation in
\emph{open} systems, with the observation that dephasing reduces quantum
fragmentation to classical. The nearest published neighbour to this project.
\item \textbf{Iadecola} \cite{Iadecola2026} --- fragmentation combined with
symmetry gives exponentially many logical qubits in degenerate pairs. The
cleanest statement of why a non-Abelian commutant is the thing to want.
\end{itemize}

\subsection{Error correction, algebraically}

\begin{itemize}
\item \textbf{Knill \& Laflamme} \cite{KL1997} --- the correctability
conditions.
\item \textbf{Zanardi \& Rasetti} \cite{ZR1997}, \textbf{Lidar, Chuang \&
Whaley} \cite{LCW1998} --- decoherence-free subspaces.
\item \textbf{Knill, Laflamme \& Viola} \cite{KLV2000} --- noiseless
\emph{subsystems}; the paper that made move~2 of \S\ref{sec:pivot} explicit.
\item \textbf{Lidar \& Whaley} \cite{LW2003} --- the review. Read this rather
than the primary papers for the picture in one place.
\item \textbf{Kribs, Laflamme \& Poulin} \cite{KLP2005} --- operator-algebra
quantum error correction, unifying DFS, noiseless subsystems and standard codes
as special cases of one algebraic statement.
\item \textbf{Blume-Kohout, Ng, Poulin \& Viola} \cite{BKNPV2008} --- every
information-preserving structure is a matrix algebra. The theorem that lets R22
\S8 say the counts \emph{are} code parameters rather than analogies.
\end{itemize}

\section{The three structure theorems}
\label{sec:theorems}

All three are stated in the finite-dimensional form, the only form this project
needs.

\subsection{Von Neumann's bicommutant theorem}

For $S \subseteq \B(\HH)$ write $S' = \{A : [A,S]=0\}$, the \emph{commutant}.

\begin{theorem}[Bicommutant]
Let $\M \subseteq \B(\HH)$ be a $*$-closed subalgebra containing $\mathbb{1}$.
Then $\M'' = \M$ in finite dimensions; in general $\M''$ is the closure of $\M$
in the weak operator topology.
\end{theorem}

\noindent Physical content: \emph{taking the commutant twice returns you to the
start}, so an algebra of observables and its commutant determine each other.
This is why one may work with either the symmetry algebra or the dynamics
algebra and lose nothing. See \cite{BR1, Blackadar2006}; for the
finite-dimensional statement, \cite{Wolf2012, Watrous2018}.

\subsection{The structure theorem (Wedderburn--Artin)}
\label{sec:wedderburn}

\begin{theorem}[Structure of a finite-dimensional $*$-algebra]
\label{thm:structure}
Let $\A \subseteq \B(\HH)$ be a unital $*$-closed subalgebra with
$\dim\HH < \infty$. Then there is a unitary under which
\[
  \HH \;\cong\; \bigoplus_k \HH_{n_k} \otimes \HH_{m_k},
  \qquad
  \A \;\cong\; \bigoplus_k M_{n_k} \otimes \mathbb{1}_{m_k},
  \qquad
  \A' \;\cong\; \bigoplus_k \mathbb{1}_{n_k} \otimes M_{m_k}.
\]
Hence $\dim\A = \sum_k n_k^2$ and $\dim\A' = \sum_k m_k^2$.
\end{theorem}

\noindent This is the theorem behind move~2 of \S\ref{sec:pivot}. The three-way
reading of the labels is worth internalising:

\begin{center}
\begin{tabular}{lll}
\toprule
symbol & superselection language & code language \\
\midrule
$k$    & sector, charge        & the classical label \\
$n_k$  & multiplicity          & the protected \emph{subsystem} (qubits) \\
$m_k$  & sector dimension      & the gauge/syndrome factor (erased) \\
\bottomrule
\end{tabular}
\end{center}

\noindent An algebra is \emph{Abelian} iff every $n_k = 1$, and a \emph{factor}
iff there is one $k$. The centre $\A \cap \A'$ is spanned by the block
identities, so $\dim(\A\cap\A')$ counts the sectors. In this project's terms,
the number of blocks is the number of superselection sectors and $\log_2 n_k$
is the number of protected qubits in block $k$.

\subsection{The asymptotic decomposition of a channel}

\begin{theorem}[Baumgartner--Narnhofer; Albert--Jiang]
\label{thm:asym}
Let $\Phi$ be a channel on a finite-dimensional $\HH$. Then
$\HH = \big(\bigoplus_k \HH_{n_k}\otimes\HH_{m_k}\big) \oplus \HH_{\rm tr}$,
and every state converges under Ces\`aro averaging to
\[
  \rho(\infty) \;=\; \bigoplus_k \; p_k \; \rho_k \otimes \sigma_k ,
\]
where $p_k = \Tr(J_k \rho_0)$ is fixed by the initial state, each $\sigma_k$ is
a \emph{fixed} full-rank state determined by $\Phi$ alone, and each $\rho_k$ is
an arbitrary state on $\HH_{n_k}$ evolving unitarily,
$\rho_k \mapsto U_k \rho_k U_k^\dagger$.
\end{theorem}

\noindent Read it as: $p_k$ is a classical label conserved from $t = 0$;
$\rho_k$ is quantum memory that survives forever; $\sigma_k$ is junk. The
transient part $\HH_{\rm tr}$ has no closed-system counterpart and is where most
of this project's subtleties live. See \cite{BN2008, AJ2014, Albert2018,
CP2016}; the code-theoretic reading is \cite{BKNPV2008}.

\section{Glossary}
\label{sec:glossary}

Grouped by theme, alphabetical within each group. Each entry gives a
definition, an intuition, and what to cite.

\subsection{Algebra}

\term{Bond algebra} The algebra generated by the local terms of a Hamiltonian
or circuit, $\A = \langle h_1, h_2, \dots\rangle$. In the fragmentation
literature the dynamics is specified by its bond algebra and the symmetries by
its commutant. \cite{MM2022}

\term{Centre} $Z(\A) = \A \cap \A'$. By Thm.~\ref{thm:structure} it is spanned
by the block identities, so $\dim Z(\A)$ is the number of superselection
sectors. ``Splitting an algebra by its centre'' means finding those blocks.
\cite{BR1}

\term{Commutant} $S' = \{A \in \B(\HH) : [A,S] = 0\ \forall S \in S\}$. Always
an algebra, even when $S$ is not; a von Neumann algebra when $S$ is $*$-closed.
\cite{BR1, MM2022}

\term{Conditional expectation} A completely positive unital
$E : \A \to \mathcal{B} \subseteq \A$ with $E|_{\mathcal{B}} = \mathrm{id}$ and
$E(bab') = b\,E(a)\,b'$ for $b,b' \in \mathcal{B}$. The algebraic version of
``trace out'' or ``project onto the conserved part''.
\cite{Wolf2012, Landsman2017}

\term{$C^*$-algebra} A Banach $*$-algebra with
$\lVert a^*a\rVert = \lVert a\rVert^2$. In finite dimensions, exactly a direct
sum of matrix algebras. \cite{Murphy1990, Landsman2017}

\term{Factor} A von Neumann algebra with trivial centre; equivalently a single
block in Thm.~\ref{thm:structure}. \cite{BR1}

\term{GNS construction} From a state $\omega$ on $\A$, builds a Hilbert space,
a representation $\pi_\omega$ and a cyclic vector $\Omega$ with
$\omega(a) = \langle\Omega, \pi_\omega(a)\Omega\rangle$. The bridge from the
algebraic formulation back to the familiar one. \cite{BR1, Strocchi2008}

\term{Operator system} A self-adjoint subspace containing $\mathbb{1}$, closed
under linear combination and adjoint but \emph{not} necessarily under products.
The correct description of the conserved quantities of an irreversible channel.
\cite{Wolf2012}

\term{Superselection sector} A block $k$ in Thm.~\ref{thm:structure}: coherent
superpositions across distinct $k$ are unobservable and dynamically irrelevant.
\cite{Haag1996, BR1}

\term{von Neumann algebra} A $*$-subalgebra $\M \subseteq \B(\HH)$ with
$\mathbb{1} \in \M$ and $\M = \M''$. In finite dimensions, any unital
$*$-subalgebra. \cite{BR1, Blackadar2006}

\subsection{Channels}

\term{Choi--Jamio\l{}kowski isomorphism} The bijection between maps
$\Phi : \B(\HH_A) \to \B(\HH_B)$ and operators
$J(\Phi) = (\Phi\otimes\mathrm{id})(|\Omega\rangle\langle\Omega|)$; $\Phi$ is
completely positive iff $J(\Phi) \ge 0$. \cite{Wolf2012, Watrous2018}

\term{Completely positive (CP)} $\Phi \otimes \mathrm{id}_n$ is positive for
every $n$. Strictly stronger than positivity, and the correct notion for a
physical map. \cite{Wolf2012}

\term{Dual (adjoint) channel} $\Phid$, defined by
$\Tr(\Phid(A)\rho) = \Tr(A\,\Phi(\rho))$ --- the Heisenberg picture. $\Phi$ is
trace preserving iff $\Phid$ is unital. Almost every algebraic statement is
made for $\Phid$. \cite{Wolf2012}

\term{Faithful state} A state of full rank; equivalently
$\omega(a^*a) = 0 \Rightarrow a = 0$. The hypothesis under which $\F$ becomes an
algebra. \cite{Frigerio1978}

\term{Fixed-point set} $\F(\Phi) = \{A : \Phi(A) = A\}$. An algebra when a
faithful invariant state exists, an operator system otherwise.
\cite{Frigerio1978, Wolf2012}

\term{Irreducible (ergodic) channel} One with no non-trivial invariant
subspace; equivalently, a unique and faithful stationary state.
\cite{CP2016, Wolf2012}

\term{Kraus operators} $\Phi(\rho) = \sum_j K_j \rho K_j^\dagger$ with
$\sum_j K_j^\dagger K_j = \mathbb{1}$. Not unique --- different sets are related
by an isometry --- which matters: a statement should be checked to be
representation-independent. \cite{Wolf2012, NC2010}

\term{Lindbladian (GKLS generator)} The generator of a continuous-time
Markovian semigroup,
$\dot\rho = -i[H,\rho] + \sum_k \big(L_k\rho L_k^\dagger -
\tfrac12\{L_k^\dagger L_k,\rho\}\big)$. \cite{BP2002, RH2012, Albert2018}

\term{Multiplicative domain} For a unital CP map $\Phi$,
$\M_\Phi = \{a : \Phi(a^*a) = \Phi(a)^*\Phi(a) \text{ and }
\Phi(aa^*) = \Phi(a)\Phi(a)^*\}$: a $C^*$-algebra on which $\Phi$ is
multiplicative. Equality in the Kadison--Schwarz inequality is the test. This is
why conserved \emph{projections} are better behaved than conserved operators
generally. \cite{Choi1974, Wolf2012}

\term{Peripheral spectrum} The eigenvalues of $\Phi$ with $|\lambda| = 1$. The
corresponding eigenoperators span the asymptotic algebra; the fixed points are
only the $\lambda = 1$ part, which is strictly smaller whenever the asymptotic
dynamics rotates. \cite{Wolf2012, Albert2018}

\term{Stinespring dilation} Every CP map is
$\Phi(\rho) = \Tr_E\!\big(V\rho V^\dagger\big)$ for an isometry $V$ into a
larger space: every channel is a unitary on a bigger system.
\cite{Wolf2012, Watrous2018}

\term{Unital channel} $\Phi(\mathbb{1}) = \mathbb{1}$; equivalently the
maximally mixed state is stationary. Unitary evolution is the special case.
R22's Prop.~1 turns on this hypothesis rather than on unitarity. \cite{Wolf2012}

\subsection{Open-system structure}

\term{Asymptotic algebra} The algebra of peripheral eigenoperators of $\Phid$
restricted to the asymptotic support; equal to
$\bigoplus_k M_{n_k}\otimes\mathbb{1}_{m_k}$. What the system remembers forever.
\cite{Albert2018, BKNPV2008}

\term{Attractivity} A subspace is \emph{invariant} if states in it stay, and
\emph{attractive} if states outside are drawn to it. These are different, and
proving the first does not give the second. \cite{TV2008}

\term{Decoherence-free subspace (DFS)} A subspace on which the channel acts
unitarily; the special case $m_k = 1$ of Thm.~\ref{thm:asym}.
\cite{ZR1997, LCW1998, LW2003}

\term{Enclosure} A subspace $\HH_0$ with
$\mathrm{supp}\,\rho \subseteq \HH_0 \Rightarrow
\mathrm{supp}\,\Phi(\rho) \subseteq \HH_0$ --- once in, never out. The
Schr\"odinger-picture invariant subspace. \cite{BN2008}

\term{Noiseless subsystem} A tensor \emph{factor} $\HH_{n_k}$ inside a block,
on which the dynamics is unitary while the co-factor $\HH_{m_k}$ is scrambled.
Strictly more general than a DFS, and the reason move~2 of \S\ref{sec:pivot}
matters. \cite{KLV2000, LW2003, KLP2005}

\term{Recurrent / transient} A state is transient if the dynamics eventually
leaves it with probability one. The recurrent part carries the asymptotics; the
transient part is invisible to Thm.~\ref{thm:asym} and is exactly what makes
$\F \ne \C$. \cite{CP2016}

\term{Strong symmetry} A unitary $U$ with $[U, K_j] = 0$ for \emph{every} Kraus
operator (for a Lindbladian, $[U,H] = [U,L_k] = 0$). Implies a charge
measurable at $t = 0$ and a block decomposition of the steady-state space.
\cite{Buca2012, AJ2014}

\term{Weak symmetry} $\Phi(U\rho U^\dagger) = U\Phi(\rho)U^\dagger$: the
generator as a whole commutes, but the individual jumps need not. Strictly
weaker, and it does \emph{not} give a superselection structure.
\cite{Buca2012, Lieu2020}

\subsection{Fragmentation}

\term{Commutant algebra (fragmentation sense)} The commutant of the bond
algebra. \textbf{Note the collision}: in the open-system literature
``commutant'' usually means the commutant of the Kraus set. \cite{MM2022}

\term{Hilbert-space fragmentation (HSF)} Exponential-in-system-size growth of
the number of dynamically disconnected sectors, not accounted for by any
conventional symmetry. Algebraically: exponential growth of the commutant's
dimension. \cite{Sala2020, Khemani2020, MM2022, MBR2022}

\term{Krylov sector} The orbit of a basis state under the dynamics; the blocks
in which fragmentation is usually counted. \cite{Sala2020, MBR2022}

\term{Strong vs.\ weak fragmentation} Whether the \emph{largest} sector occupies
a vanishing fraction of the Hilbert space (strong) or a finite one (weak).
\textbf{Unrelated} to strong/weak \emph{symmetry}. \cite{Sala2020, MBR2022}

\subsection{Error correction}

\term{Information-preserving structure} Blume-Kohout \emph{et al.}'s
classification: anything preserved by a process is a matrix algebra, isometric
to the fixed points of an associated unital process. The bridge between
Thm.~\ref{thm:asym} and code parameters. \cite{BKNPV2008}

\term{Knill--Laflamme conditions} A code with projector $P$ is correctable for
errors $\{E_a\}$ iff $P E_a^\dagger E_b P = c_{ab} P$. Detection of a single
error is the same condition with one error. \cite{KL1997, NC2010}

\term{Operator-algebra QEC (OAQEC)} The unification of DFS, noiseless subsystems
and standard codes: correctable codes are operator algebras. \cite{KLP2005}

\term{Zero-error capacity} How much can be recovered with probability exactly
one, as against on average. The distinction R22 \S3.5 needs: multistability can
have ordinary capacity while having none of this. \cite{BKNPV2008}

\section{Notational collisions}
\label{sec:collisions}

These have actually caused errors, here or in the literature.

\begin{enumerate}
\item \textbf{``Strong/weak'' means three different things.}
(i) Strong/weak \emph{symmetry}: whether each Kraus operator commutes, or only
the channel \cite{Buca2012}. (ii) Strong/weak \emph{fragmentation}: whether the
largest Krylov sector is a vanishing fraction of the space \cite{Sala2020}.
(iii) Strong-to-weak \emph{spontaneous symmetry breaking} in mixed states
\cite{Lessa2024}. Mutually independent; always say which axis.

\item \textbf{``Commutant'' has two referents.} Commutant of the bond algebra
(fragmentation, \cite{MM2022}) versus commutant of the Kraus set (open systems,
\cite{Buca2012}). They coincide for closed dynamics and not otherwise.

\item \textbf{``Subsystem'' has two referents.} A tensor factor of the whole
space (physics-degree sense) versus a factor $\HH_{n_k}$ of one block
(algebraic sense, \cite{KLV2000, KLP2005}). Nearly every sentence in the QEC
literature means the second.

\item \textbf{Conserved quantity $\ne$ symmetry.} $\C \subseteq \F$ always, with
equality iff a faithful invariant state exists \cite{Frigerio1978}. Writing
``conserved quantity'' where ``strong symmetry'' is meant is a substantive
error in open systems, not loose phrasing.

\item \textbf{Fixed points $\ne$ asymptotic algebra.} The fixed points of
$\Phid$ are the $\lambda = 1$ eigenoperators; the asymptotic algebra is the
whole peripheral spectrum. They differ whenever the asymptotic dynamics
rotates --- for a unitary channel, dramatically.
\end{enumerate}

\section{Where each piece is used in this project}
\label{sec:project}

\begin{center}
\small
\begin{tabular}{lll}
\toprule
concept & report & what is computed \\
\midrule
commutant algebra, HSF criterion & R21, R22 &
  $n_{\mathrm{wcc}} = \dim(\C \cap \D)$ \\
strong symmetry (Bu\v{c}a--Prosen) & R22 &
  weak components of the support graph \\
enclosure (Baumgartner--Narnhofer) & R22, R24 &
  terminal strong components \\
conserved quantities $\F$ & R22, R24 &
  absorption probabilities; the squeeze \\
multiplicative domain, projections & R22, R24 &
  every conserved projection is a strong symmetry \\
Thm.~\ref{thm:asym} decomposition & R22, R24 &
  $\bigoplus_k M_{n_k}\otimes\sigma_{m_k}$, exactly \\
noiseless subsystem, DFS & R23, R24 &
  $m_k = 1$ for all eight rules \\
Knill--Laflamme, code distance & R24 &
  distance $1$; no subcode helps \\
information-preserving structure & R22 &
  bits vs.\ qubits vs.\ nothing \\
\bottomrule
\end{tabular}
\end{center}

\noindent Two results of this project are statements \emph{about} the framework
rather than applications of it, and are the natural places for a thesis to
claim novelty. First, R22's theorem that the projections in $\F$ are exactly
the indicators of unions of weak components, so the Boolean algebra of
always-answerable questions cannot grow with the number of attractors. Second,
R24's correction that $\dim(\F\cap\D) = n_{\mathrm{rec}}$ holds only for the
monitored chain, the coherent statement being the squeeze
$n_{\mathrm{wcc}} \le \dim(\F\cap\D) \le n_{\mathrm{rec}}$, with $W_{29}$ at
$N=4$ as an explicit counterexample to equality.

\section*{Bibliography note}
\label{sec:biblio}
\addcontentsline{toc}{section}{Bibliography note}

Entries marked $\checkmark$ were verified against the publisher or arXiv while
this document was written. Entries marked $\dagger$ are standard textbooks,
cited from canonical bibliographic data and not re-verified; check the edition
before citing a page number.

\begin{thebibliography}{99}

\bibitem{AJ2014} $\checkmark$ V.~V.~Albert and L.~Jiang,
  \emph{Symmetries and conserved quantities in Lindblad master equations},
  Phys.\ Rev.\ A \textbf{89}, 022118 (2014); arXiv:1310.1523.

\bibitem{Albert2018} $\checkmark$ V.~V.~Albert,
  \emph{Lindbladians with multiple steady states: theory and applications},
  arXiv:1802.00010 (2018). 134 pp.; a thesis that reads as a review.

\bibitem{BN2008} $\checkmark$ B.~Baumgartner and H.~Narnhofer,
  \emph{Analysis of quantum semigroups with GKS--Lindblad generators II:
  General},
  J.\ Phys.\ A \textbf{41}, 395303 (2008); arXiv:0806.3164.

\bibitem{Blackadar2006} $\dagger$ B.~Blackadar,
  \emph{Operator Algebras: Theory of $C^*$-Algebras and von Neumann Algebras},
  Springer (2006).

\bibitem{BKNPV2008} $\checkmark$ R.~Blume-Kohout, H.~K.~Ng, D.~Poulin and
  L.~Viola,
  \emph{Characterizing the structure of preserved information in quantum
  processes},
  Phys.\ Rev.\ Lett.\ \textbf{100}, 030501 (2008).

\bibitem{BR1} $\dagger$ O.~Bratteli and D.~W.~Robinson,
  \emph{Operator Algebras and Quantum Statistical Mechanics 1: $C^*$- and
  $W^*$-Algebras, Symmetry Groups, Decomposition of States},
  2nd ed., Springer (1987).

\bibitem{BP2002} $\dagger$ H.-P.~Breuer and F.~Petruccione,
  \emph{The Theory of Open Quantum Systems}, Oxford University Press (2002).

\bibitem{Buca2012} $\checkmark$ B.~Bu\v{c}a and T.~Prosen,
  \emph{A note on symmetry reductions of the Lindblad equation: transport in
  constrained open spin chains},
  New J.\ Phys.\ \textbf{14}, 073007 (2012); arXiv:1203.0943.

\bibitem{CP2016} $\checkmark$ R.~Carbone and Y.~Pautrat,
  \emph{Irreducible decompositions and stationary states of quantum channels},
  Rep.\ Math.\ Phys.\ \textbf{77}, 293 (2016); arXiv:1507.08404.

\bibitem{Choi1974} $\checkmark$ M.-D.~Choi,
  \emph{A Schwarz inequality for positive linear maps on $C^*$-algebras},
  Illinois J.\ Math.\ \textbf{18}, 565 (1974). The multiplicative domain.

\bibitem{Frigerio1978} $\checkmark$ A.~Frigerio,
  \emph{Stationary states of quantum dynamical semigroups},
  Commun.\ Math.\ Phys.\ \textbf{63}, 269 (1978).

\bibitem{Haag1996} $\dagger$ R.~Haag,
  \emph{Local Quantum Physics: Fields, Particles, Algebras},
  2nd ed., Springer (1996).

\bibitem{HZ2011} $\dagger$ T.~Heinosaari and M.~Ziman,
  \emph{The Mathematical Language of Quantum Theory}, Cambridge (2011).

\bibitem{Iadecola2026} $\checkmark$ T.~Iadecola,
  \emph{Symmetry fragmentation},
  Phys.\ Rev.\ Lett.\ \textbf{136}, 100401 (2026); arXiv:2510.06333.

\bibitem{Khemani2020} $\checkmark$ V.~Khemani, M.~Hermele and R.~Nandkishore,
  \emph{Localization from Hilbert space shattering: from theory to physical
  realizations},
  Phys.\ Rev.\ B \textbf{101}, 174204 (2020); arXiv:1910.01137.

\bibitem{KL1997} $\checkmark$ E.~Knill and R.~Laflamme,
  \emph{Theory of quantum error-correcting codes},
  Phys.\ Rev.\ A \textbf{55}, 900 (1997).

\bibitem{KLV2000} $\checkmark$ E.~Knill, R.~Laflamme and L.~Viola,
  \emph{Theory of quantum error correction for general noise},
  Phys.\ Rev.\ Lett.\ \textbf{84}, 2525 (2000).

\bibitem{KLP2005} $\checkmark$ D.~W.~Kribs, R.~Laflamme and D.~Poulin,
  \emph{A unified and generalized approach to quantum error correction},
  Phys.\ Rev.\ Lett.\ \textbf{94}, 180501 (2005).

\bibitem{Landsman2017} $\checkmark$ K.~Landsman,
  \emph{Foundations of Quantum Theory: From Classical Concepts to Operator
  Algebras},
  Fundamental Theories of Physics \textbf{188}, Springer (2017).
  \textbf{Open access.}

\bibitem{Lessa2024} $\checkmark$ L.~A.~Lessa, R.~Ma, J.-H.~Zhang, Z.~Bi,
  M.~Cheng and C.~Wang,
  \emph{Strong-to-weak spontaneous symmetry breaking in mixed quantum states},
  PRX Quantum \textbf{6}, 010344; arXiv:2405.03639.

\bibitem{Li2023} $\checkmark$ Y.~Li, P.~Sala and F.~Pollmann,
  \emph{Hilbert space fragmentation in open quantum systems},
  Phys.\ Rev.\ Research \textbf{5}, 043239 (2023); arXiv:2305.06918.

\bibitem{LCW1998} $\checkmark$ D.~A.~Lidar, I.~L.~Chuang and K.~B.~Whaley,
  \emph{Decoherence-free subspaces for quantum computation},
  Phys.\ Rev.\ Lett.\ \textbf{81}, 2594 (1998).

\bibitem{LW2003} $\checkmark$ D.~A.~Lidar and K.~B.~Whaley,
  \emph{Decoherence-free subspaces and subsystems}, in
  \emph{Irreversible Quantum Dynamics}, eds.\ F.~Benatti and R.~Floreanini,
  Springer Lecture Notes in Physics \textbf{622}, 83--120 (2003);
  arXiv:quant-ph/0301032. \textbf{The review to read.}

\bibitem{Lieu2020} $\checkmark$ S.~Lieu, R.~Belyansky, J.~T.~Young,
  R.~Lundgren, V.~V.~Albert and A.~V.~Gorshkov,
  \emph{Symmetry breaking and error correction in open quantum systems},
  Phys.\ Rev.\ Lett.\ \textbf{125}, 240405 (2020); arXiv:2008.02816.

\bibitem{MM2022} $\checkmark$ S.~Moudgalya and O.~I.~Motrunich,
  \emph{Hilbert space fragmentation and commutant algebras},
  Phys.\ Rev.\ X \textbf{12}, 011050 (2022); arXiv:2108.10324.

\bibitem{MM2024} $\checkmark$ S.~Moudgalya and O.~I.~Motrunich,
  \emph{Symmetries as ground states of local superoperators: hydrodynamic
  implications},
  PRX Quantum \textbf{5}, 040330 (2024); arXiv:2309.15167.

\bibitem{MBR2022} $\checkmark$ S.~Moudgalya, B.~A.~Bernevig and N.~Regnault,
  \emph{Quantum many-body scars and Hilbert space fragmentation: a review of
  exact results},
  Rep.\ Prog.\ Phys.\ \textbf{85}, 086501 (2022); arXiv:2109.00548.
  \textbf{The review to cite.}

\bibitem{Murphy1990} $\dagger$ G.~J.~Murphy,
  \emph{$C^*$-Algebras and Operator Theory}, Academic Press (1990).

\bibitem{NC2010} $\dagger$ M.~A.~Nielsen and I.~L.~Chuang,
  \emph{Quantum Computation and Quantum Information},
  10th anniversary ed., Cambridge (2010).

\bibitem{RH2012} $\dagger$ \'A.~Rivas and S.~F.~Huelga,
  \emph{Open Quantum Systems: An Introduction}, SpringerBriefs (2012).

\bibitem{Sala2020} $\checkmark$ P.~Sala, T.~Rakovszky, R.~Verresen, M.~Knap and
  F.~Pollmann,
  \emph{Ergodicity breaking arising from Hilbert space fragmentation in
  dipole-conserving Hamiltonians},
  Phys.\ Rev.\ X \textbf{10}, 011047 (2020); arXiv:1904.04266.

\bibitem{Strocchi2008} $\dagger$ F.~Strocchi,
  \emph{An Introduction to the Mathematical Structure of Quantum Mechanics},
  2nd ed., World Scientific (2008). \textbf{Gentlest entry point.}

\bibitem{Takesaki2002} $\dagger$ M.~Takesaki,
  \emph{Theory of Operator Algebras I--III}, Springer (2002--2003).

\bibitem{TV2008} $\checkmark$ F.~Ticozzi and L.~Viola,
  \emph{Quantum Markovian subsystems: invariance, attractivity, and control},
  IEEE Trans.\ Autom.\ Control \textbf{53}, 2048 (2008); arXiv:0705.1372.

\bibitem{Watrous2018} $\dagger$ J.~Watrous,
  \emph{The Theory of Quantum Information}, Cambridge (2018).

\bibitem{Wolf2012} $\checkmark$ M.~M.~Wolf,
  \emph{Quantum Channels \& Operations: Guided Tour},
  lecture notes, 5 July 2012. \textbf{Start here.}\\
  \href{https://mediatum.ub.tum.de/doc/1701036/1701036.pdf}%
  {mediatum.ub.tum.de/doc/1701036/1701036.pdf}

\bibitem{ZR1997} $\checkmark$ P.~Zanardi and M.~Rasetti,
  \emph{Noiseless quantum codes},
  Phys.\ Rev.\ Lett.\ \textbf{79}, 3306 (1997).

\end{thebibliography}

\end{document}
